\documentclass[10pt,a4paper]{amsart}
\usepackage[margin=19mm]{geometry}
\usepackage{amsmath,amssymb,mathtools}
\usepackage[scr=rsfs,cal=euler]{mathalfa}
\usepackage{xfrac}
\usepackage[unicode,hidelinks,bookmarks=false]{hyperref}
\newcommand{\ie}{i.e.\ }
\newcommand{\cf}{cf.\ }
\newcommand{\resp}{respectively\ }
\newcommand{\Subsection}{\subsection}
\providecommand{\nobreakdash}{\nobreak\hbox{-}}
\setlength{\emergencystretch}{2em}
\multlinegap0pt
\usepackage{bm}
\usepackage{bbm}
\usepackage{dsfont}
\usepackage{xspace}
\usepackage{cancel}
\usepackage{mathtools}
\usepackage{amsthm}
\makeatletter
\@ifundefined{theorem}{\newtheorem{theorem}{Theorem}}{}
\@ifundefined{observation}{\newtheorem{observation}[theorem]{Observation}}{}
\@ifundefined{proposition}{\newtheorem{proposition}[theorem]{Proposition}}{}
\makeatother
\makeatletter
\DeclareMathOperator{\diag}{diag}
\expandafter\ifx\csname frontmatter\endcsname\relax
\newenvironment{frontmatter}{}{}
\fi
\expandafter\ifx\csname keyword\endcsname\relax
\newenvironment{keyword}{\small \textbf{Keywords:}}{}
\fi
\newcommand{\myref}[1]{\textnormal{(\hyperlink{#1}{#1}})}
\newcommand{\myrefmath}[1]{(\hyperlink{#1}{\mathrm{#1}})}
\newcommand{\mytag}[1]{(\hypertarget{#1}{\mathrm{#1}})}
\expandafter\ifx\csname subalgorithms\endcsname\relax
\newenvironment{subalgorithms}{%
  \refstepcounter{algorithm}%
  \protected@edef\theparentalgorithm{\thealgorithm}%
  \setcounter{parentalgorithm}{\value{algorithm}}%
  \setcounter{algorithm}{0}%
  \def\thealgorithm{\theparentalgorithm\alph{algorithm}}%
  \ignorespaces
}{%
  \setcounter{algorithm}{\value{parentalgorithm}}%
  \ignorespacesafterend
}
\fi
\expandafter\ifx\csname varsubequations\endcsname\relax
\newenvironment{varsubequations}[1]
 {%
  \addtocounter{equation}{-1}%
  \begin{subequations}
  \renewcommand{\theparentequation}{#1.}%
  \def\@currentlabel{#1}%
 }
\fi
\makeatother
\title[An semidefinite programming-based $\varepsilon$-constraint m]{An semidefinite programming-based $\varepsilon$-constraint method for the bi-objective single-row facility layout problem}
\author{Christof Brandstetter, Elisabeth Gaar, Markus Sinnl}
\date{Source: arXiv:2607.00611v1 (CC BY 4.0)}
\begin{document}
\maketitle

\noindent\emph{Source and licence.} An semidefinite programming-based $\varepsilon$-constraint method for the bi-objective single-row facility layout problem. Version arXiv:2607.00611v1 (CC BY 4.0), distributed under the Creative Commons Attribution 4.0 International licence, as recorded in the arXiv licence metadata for this exact version and shown on the abstract page https://arxiv.org/abs/2607.00611v1. Full rights notice: licenses/arxiv-2607.00611-CC-BY-4.0.txt.\par\smallskip

\noindent\textit{Editorial scope.} Counted unit: the Proposition whose source label is \texttt{prop:red} in the article (source0.tex lines 692--724, source label \texttt{prop:red}), delivered together with the article's own complete original proof of it (source0.tex lines 726--774, one \texttt{proof} environment of 49 lines). No mathematics is written, repaired or re-derived: statement and proof are the article's own text line for line. Required context retained uncounted: eq:transformationileqj (lines 364--367); eq:MO-SRFLP-3Cycle (lines 392--401); eq:MO-SRFLP-diag (lines 392--401); eq:MO-SRFLP-psd (lines 396--399); eq:BOR-down (lines 492--498); eq:BOR-eps (lines 492--498); eq:BOR-obj (lines 492--498); eq:BOR-rest (lines 492--498); obs:red (lines 671--678). Every definition, display and label of those lines is retained unchanged. Omitted from this package and not redistributed: 1--363, 368--391, 400--491, 499--670, 679--691, 725--725, 775--1410. Nothing inside a retained excerpt is omitted or altered: every retained body is delivered exactly as the article's lines are written, so no omission receipt of any kind accompanies this package. Citations: the retained excerpts carry no citation command at all, so no bibliography entry of the article is transcribed and no reference is redistributed. Reference adaptation: none was needed and none was made; every reference inside the delivered unit resolves inside it, so no unresolved reference or citation is produced. Printed slips kept literal: no printed word, symbol, hypothesis or number of the counted statement or of its proof was altered. Layout and class: the article's own class is \texttt{article} with packages including enumitem, algorithm, algpseudocode, inputenc, fontenc, amsmath, amssymb, mathtools, amsthm, amsfonts, graphicx, array, hyperref, float; the wrapper is the lane's pinned \texttt{amsart} editorial wrapper at 10pt on A4, which restates the theorem-like environments the retained text needs and adds the article's own macro definitions that the retained text needs (\texttt{\textbackslash diag}, \texttt{\textbackslash myref}, \texttt{\textbackslash myrefmath}, \texttt{\textbackslash mytag}), with extra wrapper packages (bm, bbm, dsfont, xspace, cancel, mathtools). The delivered numbering is the unit's own. Assets: no retained span contains a figure environment, a graphics-inclusion command, a PDF-page inclusion, a table or a TikZ picture; no image file is copied into this package, the package declares no asset dependency and no image file is redistributed with it. Licence: version arXiv:2607.00611v1 of this article is distributed under the Creative Commons Attribution 4.0 International licence, as recorded in the arXiv licence metadata for this exact version and shown on the abstract page https://arxiv.org/abs/2607.00611v1. A full rights notice is delivered at licenses/arxiv-2607.00611-CC-BY-4.0.txt. Characters: every retained excerpt is pure ASCII, so the delivered engine, XeLaTeX with its default Latin Modern text font, reports no missing character.
\begin{align}\label{eq:transformationileqj}
    X_{ij,kh} = \tau(i,j) \tau(k,h) X_{\min\{i,j\}\max\{i,j\}, \min\{k,h\}\max\{k,h\}} && 
    \forall i, j,k,h \in I 
\end{align}

    \begin{align}
        \min        \quad   & (K_1+\langle C_1, X\rangle,\dots,K_p+\langle C_p, X\rangle) \\
        \text{s.t.} \quad  & \sum_{\substack{k \in I, \\ i \neq k \neq j}} X_{ij,jk} - X_{ij,ik} - X_{ik,jk} = -(n-2) &\qquad& \forall i, j \in I, i < j \label{eq:MO-SRFLP-3Cycle}\\
                            & \diag(X) = e \label{eq:MO-SRFLP-diag}\\
                            & \begin{pmatrix}
                                1 & x^T \\
                                x & X
                            \end{pmatrix}\succeq 0 \label{eq:MO-SRFLP-psd}\\
                             & x_{ij} \in \{-1,1\} && \forall i, j \in I, i < j. \label{eq:MO-SRFLP-xint} 
    \end{align} 

    \begin{align}
        \mytag{SDP(\varepsilon,I^-,I^+)} \quad z = \min \quad &K_1 +\langle C_1,X \rangle \label{eq:BOR-obj}\\
        \text{s.t.} \quad & K_2 + \langle C_2, X \rangle \leq \varepsilon \label{eq:BOR-eps} \\ 
        & \eqref{eq:MO-SRFLP-3Cycle}-\eqref{eq:MO-SRFLP-psd} \\
        &x_{ij}=-1 &\quad&\forall ij \in I^- \label{eq:BOR-down} \\
        &x_{ij}=1 &&\forall ij \in I^+ \label{eq:BOR-rest}.
    \end{align}

\begin{observation}
\label{obs:red}
In \myref{SDP($\varepsilon,I^-,I^+$)}, any constraint~\eqref{eq:BOR-down} and~\eqref{eq:BOR-rest}, i.e., any constraint of the form $x_{ij} = \sigma$ for some $ij \in I^2$ with $\sigma \in \{-1,1\}$, implies 
\begin{align*}
    X_{ij,kh} = X_{kh,ij} = \sigma x_{kh} \quad \forall kh \in I^2.
\end{align*} 
Thus, the column of $X$ indexed by $ij$ equals $\sigma x$ and the row of $X$ indexed by $ij$ equals $\sigma x^T$.
\end{observation}

\begin{proposition}    \label{prop:red}
Consider a node in the B\&B tree with given $\varepsilon$, $I^-$ and $I^+$.
Let $I^2_F = I^- \cup I^+$ and $I^2_R = I^2\setminus I^2_F$ be the fixed and non-fixed indices of $I^2$, respectively. 
Let $\sigma^F \in \{-1,1\}^{|I^2_F|}$ be the vector of all fixings in $I^2_F$. 

Then $\myrefmath{SDP(\varepsilon,I^-,I^+)}$ is equivalent to 
\begin{subequations}
    \begin{align}
        \mytag{SDP^R(\varepsilon,I^-,I^+)} \quad z = \min \quad &K_1^R +\langle C_1^R,X^R \rangle + \langle c_1^R, x^R \rangle \label{eq:SDPR-obj}\\
        \text{s.t.} \quad & K_2^R + \langle C_2^R, X^R \rangle + \langle c_2^R, x^R \rangle \leq \varepsilon \label{eq:SDPR-eps} \\ 
         & \sum_{\substack{k \in I, \\ i \neq k \neq j}} \widetilde{X}_{ij,jk} - \widetilde{X}_{ij,ik} - \widetilde{X}_{ik,jk} = -(n-2) &\quad&  \label{eq:SDPR-3Cycle}
         \forall i, j \in I, i < j \\
                            & \diag(X^R) = e \label{eq:SDPR-diag} \\
                            & \begin{pmatrix}
                                1 & (x^R)^T \\
                                x^R & X^R
                            \end{pmatrix}\succeq 0 \label{eq:SDPR-psd},
    \end{align}
\end{subequations}
where $x^R \in \mathbb R^{|I^2_R|}$, $X^R \in \mathbb R^{|I^2_R|\times|I^2_R|}$, 

\begin{align} \label{eq:fixingImplications}
\widetilde{X}_{ij,kh} = 
    \begin{cases}
        X^R_{ij,kh} & \text{if } ij,kh \in I^2_R \\
        \sigma^F_{ij}x^R_{kh} & \text{if } ij \in I^2_F, kh \in I^2_R\\        
        \sigma^F_{kh}x^R_{ij} & \text{if } kh \in I^2_F, ij \in I^2_R\\
        \sigma^F_{ij}\sigma^F_{kh} & \text{if } ij,kh \in I^2_F,
    \end{cases}
\end{align}
for all $ij, kh \in I^2$ (and thus for all $i,j,k,h \in I$ with $i\neq j$ and $k \neq h$ using~\eqref{eq:transformationileqj} to obtain $i<j$ and $k<h$), 
and $K_q^R = K_q + \langle C_{q_{I^2_F, I^2_F}}, \sigma^F(\sigma^F)^T \rangle$, $C_q^R = C_{q_{I^2_R, I^2_R}}$ and $c_q^R = 2(C_{q_{I^2_F,I^2_R}})^T \sigma^F$ for all $q\in \{1,2\}$. 
\end{proposition}

\begin{proof}
    The variable $x$ in $\myrefmath{SDP(\varepsilon,I^-,I^+)}$ can be split into $x^R \in \mathbb R^{|I^2_R|}$ and $x^F \in \mathbb R^{|I^2_F|}$ with $x^F$ representing all entries of $x$ that correspond to fixings in $I^2_F$ as 
    $x=\begin{pmatrix} x^R \\ x^F \end{pmatrix}$. 
    Similarly, $X$ can be split into 
    $X=\begin{pmatrix} X^R & (X^{FR})^T\\ X^{FR} & X^F \end{pmatrix}$ with $X^R \in \mathbb R^{|I^2_R|\times|I^2_R|}$, $X^{FR} \in \mathbb R^{|I^2_F|\times|I^2_R|}$ and $X^F \in \mathbb R^{|I^2_F|\times|I^2_F|}$.

    With this notation, the constraints~\eqref{eq:BOR-down} and~\eqref{eq:BOR-rest} of $\myrefmath{SDP(\varepsilon,I^-,I^+)}$ are equivalent to $x^F = \sigma^F$. From Observation~\ref{obs:red} it then follows that $X^{FR} = \sigma^F(x^R)^T$ and $X^F = \sigma^F(x^F)^T =\sigma^F(\sigma^F)^T$ hold.
    As a consequence, it is easy to see that constraint~\eqref{eq:MO-SRFLP-diag} is equivalent to~\eqref{eq:SDPR-diag}. Moreover, the constraint~\eqref{eq:MO-SRFLP-3Cycle} is equivalent to~\eqref{eq:SDPR-3Cycle} by using the definition of $\widetilde{X}$. 
    Furthermore, by easy computations that exploit the properties of the Frobenius inner product~\eqref{eq:BOR-obj} and~\eqref{eq:BOR-eps} and equivalent to~\eqref{eq:SDPR-obj} and~\eqref{eq:SDPR-eps}.
    
\newcommand{\myVecE}{a} %extension
\newcommand{\myVecR}{b} %remaining 
\newcommand{\myVecF}{c} %fixed
Finally, for each $\myVecE \in \mathbb{R}$, $\myVecR \in \mathbb R^{|I^2_R|}$ and $\myVecF \in \mathbb R^{|I^2_F|}$ it holds that 
\begingroup
\allowdisplaybreaks
\begin{align*}
    \begin{pmatrix} \myVecE \\ \myVecR \\ \myVecF \end{pmatrix}^T
    &\begin{pmatrix}
        1   & (x^R)^T & (x^F)^T \\
        x^R & X^R   & (X^{FR})^T \\
        x^F & X^{FR} & X^F  \end{pmatrix}
    \begin{pmatrix}\myVecE \\ \myVecR \\ \myVecF \end{pmatrix} 
    = \begin{pmatrix} \myVecE \\ \myVecR \\ \myVecF \end{pmatrix}^T 
    \begin{pmatrix*}[l]
        \myVecE + (x^R)^T \myVecR + (x^F)^T \myVecF \\ 
        \myVecE x^R + X^R \myVecR + (X^{FR})^T \myVecF\\ 
        \myVecE x^F + X^{FR} \myVecR + X^F \myVecF \end{pmatrix*}\\
    = & \myVecE^2 + \myVecE (x^R)^T \myVecR + \myVecE (x^F)^T \myVecF 
     + \myVecE \myVecR^T x^R + \myVecR^T X^R \myVecR + \myVecR^T (X^{FR})^T \myVecF 
     + \myVecE \myVecF^T x^F + \myVecF^T X^{FR} \myVecR + \myVecF^T X^F \myVecF\\
    = & \myVecE^2 + 2 \myVecE \myVecR^T x^R + 2 \myVecE \myVecF^T x^F 
      + \myVecR^T X^R \myVecR + 2 \myVecF^T X^{FR} \myVecR 
      + \myVecF^T X^F \myVecF\\ 
    = & \myVecE^2 + 2 \myVecE \myVecR^T x^R + 2 \myVecE \myVecF^T \sigma^F 
      + \myVecR^T X^R \myVecR + 2 \myVecF^T \sigma^F(x^R)^T \myVecR 
      + \myVecF^T \sigma^F(\sigma^F)^T \myVecF\\ 
    = & (\myVecE + \myVecF^T \sigma^F)^2 + 2 (\myVecE + \myVecF^T \sigma^F) \myVecR^T x^R 
      + \myVecR^T X^R \myVecR \\ 
    = & \begin{pmatrix} \myVecE + \myVecF^T \sigma^F \\ \myVecR \end{pmatrix}^T
    \begin{pmatrix}
        1 & (x^R)^T \\
        x^R & X^R \\
    \end{pmatrix} 
    \begin{pmatrix} \myVecE + \myVecF^T \sigma^F \\ \myVecR \end{pmatrix},
\end{align*}
\endgroup
implying that~\eqref{eq:MO-SRFLP-psd} is equivalent to~\eqref{eq:SDPR-psd}. Thus, $\myrefmath{SDP(\varepsilon,I^-,I^+)}$ is equivalent to $\myrefmath{SDP^R(\varepsilon,I^-,I^+)}$.
\end{proof}

\end{document}
